% !TeX root = ../教材题库_全部条目_转写工作稿.tex
% Source: XJ-B1-C03-PR-P084-Y1532
% 第3章 / 3.2 函数的基本性质 / 教材 p.84 / PDF 第91页
% OCR未校对：公式与文字均可能有误；未完成数学TeX转写；禁止直接发给学生
% 原题及相邻原解截图（核对用）：../题目原图/XJ-B1-C03-PR-P084-Y1532-01.webp
% 本文件只在转写完成后作为题干使用；OCR 条目尚不是可靠题干。
\begin{bankitem}{p.84 练习 第4题}{第3章 / 3.2 函数的基本性质 / 教材 p.84 / PDF 第91页}{XJ-B1-C03-PR-P084-Y1532}{OCR未校对：公式与文字均可能有误；未完成数学TeX转写；禁止直接发给学生}
\begin{ocrdraft}
[教材p.84 / PDF第91页]
\par
4．设g(x）是定义于（一。。，十亠）上的函数，且f(x)=g(x)\ocrglyph{00AF}g(\ocrglyph{00AF}x),讨论了(x）
\par
的奇偶性；如果在[0，十）上丆）一亠，求丆0）在（一亠，明上的解析式．
\end{ocrdraft}
\end{bankitem}
